\documentclass[12pt,a4paper]{article}
\usepackage{hojasemanal}

% ─── EDITAR AQUÍ ─────────────────────────────────────────────────────────────
\tema{Diagnóstico}
\fecha{dd/mm/aaaa}   % opcional: borrar la línea para no mostrar fecha
% ─────────────────────────────────────────────────────────────────────────────

\begin{document}

    \encabezado

    \begin{indicaciones}
        \item Si no logras terminar una demostración, escribe hasta donde llegaste.
        \item Anota en la bitácora cuánto te tomó cada ejercicio.
    \end{indicaciones}

    \begin{recordatorio}
        Un entero $n$ es \textit{par} si existe $k\in\Z$ tal que $n=2k$, e \textit{impar} si existe $k\in\Z$ tal que $n=2k+1$. Si $a,b\in\Z$, decimos que $a$ \textit{divide} a $b$, y escribimos $a\mid b$, si existe $k\in\Z$ tal que $b=ak$. Para conjuntos $A$ y $B$, $A\setminus B=\{x\in A:x\notin B\}$.
    \end{recordatorio}

    \parte{Números enteros}

    \begin{ejercicio}[1]
        Demuestra que la suma de dos enteros pares es par.
    \end{ejercicio}

    \begin{ejercicio}[1]
        Demuestra que si $n$ es un entero impar, entonces $n^2$ es impar.
    \end{ejercicio}

    \begin{ejercicio}[1]
        Sean $a,b,c\in\Z$. Demuestra que si $a\mid b$ y $b\mid c$, entonces $a\mid c$.
    \end{ejercicio}

    \begin{ejercicio}[2]
        Demuestra que si $n$ es un entero y $n^2$ es par, entonces $n$ es par.
    \end{ejercicio}

    \begin{ejercicio}[2]
        Demuestra que no existe $M\in\Z$ tal que $n\leq M$ para todo $n\in\Z$.
    \end{ejercicio}

    \parte{Conjuntos}

    \begin{ejercicio}[1]
        Sean $A$, $B$ y $C$ conjuntos con $A\subseteq B$. Demuestra que $A\cap C\subseteq B\cap C$.
    \end{ejercicio}

    \begin{ejercicio}[2]
        Sean $A$, $B$ y $C$ conjuntos. Demuestra que
        \begin{equation*}
            A\setminus(B\cup C)=(A\setminus B)\cap(A\setminus C).
        \end{equation*}
    \end{ejercicio}

    \parte{Números reales}

    \begin{ejercicio}[2]
        Demuestra que si $x$ es un número real con $x>0$, entonces
        \begin{equation*}
            x+\frac{1}{x}\geq2.
        \end{equation*}
    \end{ejercicio}

    \parte{Inducción}

    \begin{ejercicio}[2]
        Demuestra que para todo $n\in\N$,
        \begin{equation*}
            1+2+\cdots+n=\frac{n(n+1)}{2}.
        \end{equation*}
    \end{ejercicio}

    \begin{ejercicio}[2]
        Demuestra que $2^n\geq n+1$ para todo $n\in\N$.
    \end{ejercicio}

    \parte{Reto (opcional)}

    \begin{ejercicio}[3]
        Demuestra que $\sqrt{2}$ es irracional.
    \end{ejercicio}

    \begin{ejercicio}[3]
        Demuestra que hay una infinidad de números primos.
    \end{ejercicio}

    \bitacora

\end{document}
